\documentclass[10pt,a4paper]{amsart}
\usepackage{amsmath,amssymb,amsthm,graphicx}
\usepackage{verbatim}
\newenvironment{reference}{\comment}{\endcomment}
%\newenvironment{reference}{}{}
\newenvironment{slogan}{\comment}{\endcomment}
\newenvironment{history}{\comment}{\endcomment}

% For commutative diagrams we use Xy-pic
\usepackage[all]{xy}

% We use 2cell for 2-commutative diagrams.
\xyoption{2cell}
% Theorem environments.
%
\theoremstyle{plain}
\newtheorem{theorem}[subsection]{Theorem}
\newtheorem{proposition}[subsection]{Proposition}
\newtheorem{lemma}[subsection]{Lemma}

\theoremstyle{definition}
\newtheorem{definition}[subsection]{Definition}
\newtheorem{example}[subsection]{Example}
\newtheorem{exercise}[subsection]{Exercise}
\newtheorem{situation}[subsection]{Situation}

\theoremstyle{remark}
\newtheorem{remark}[subsection]{Remark}
\newtheorem{remarks}[subsection]{Remarks}

\numberwithin{equation}{subsection}

% Macros
%
\def\lim{\mathop{\mathrm{lim}}\nolimits}
\def\colim{\mathop{\mathrm{colim}}\nolimits}
\def\Spec{\mathop{\mathrm{Spec}}}
\def\Hom{\mathop{\mathrm{Hom}}\nolimits}
\def\Ext{\mathop{\mathrm{Ext}}\nolimits}
\def\SheafHom{\mathop{\mathcal{H}\!\mathit{om}}\nolimits}
\def\SheafExt{\mathop{\mathcal{E}\!\mathit{xt}}\nolimits}
\def\Sch{\mathit{Sch}}
\def\Mor{\mathop{\mathrm{Mor}}\nolimits}
\def\Ob{\mathop{\mathrm{Ob}}\nolimits}
\def\Sh{\mathop{\mathit{Sh}}\nolimits}
\def\NL{\mathop{N\!L}\nolimits}
\def\CH{\mathop{\mathrm{CH}}\nolimits}
\def\proetale{{pro\text{-}\acute{e}tale}}
\def\etale{{\acute{e}tale}}
\def\QCoh{\mathit{QCoh}}
\def\Ker{\mathop{\mathrm{Ker}}}
\def\Im{\mathop{\mathrm{Im}}}
\def\Coker{\mathop{\mathrm{Coker}}}
\def\Coim{\mathop{\mathrm{Coim}}}

% Boxtimes
%
\DeclareMathSymbol{\boxtimes}{\mathbin}{AMSa}{"02}

%
% Macros for moduli stacks/spaces
%
\def\QCohstack{\mathcal{QC}\!\mathit{oh}}
\def\Cohstack{\mathcal{C}\!\mathit{oh}}
\def\Spacesstack{\mathcal{S}\!\mathit{paces}}
\def\Quotfunctor{\mathrm{Quot}}
\def\Hilbfunctor{\mathrm{Hilb}}
\def\Curvesstack{\mathcal{C}\!\mathit{urves}}
\def\Polarizedstack{\mathcal{P}\!\mathit{olarized}}
\def\Complexesstack{\mathcal{C}\!\mathit{omplexes}}
% \Pic is the operator that assigns to X its picard group, usage \Pic(X)
% \Picardstack_{X/B} denotes the Picard stack of X over B
% \Picardfunctor_{X/B} denotes the Picard functor of X over B
\def\Pic{\mathop{\mathrm{Pic}}\nolimits}
\def\Picardstack{\mathcal{P}\!\mathit{ic}}
\def\Picardfunctor{\mathrm{Pic}}
\def\Deformationcategory{\mathcal{D}\!\mathit{ef}}
\usepackage{cleveref}
\crefname{originallemma}{original lemma}{original lemmas}
\Crefname{originallemma}{Original lemma}{Original lemmas}
\makeatletter
\newlabel{chow-lemma-pullback-base-change}{{0FVJ}{source}{}{sourcepointer-1}{}}
\newlabel{chow-lemma-pullback-base-change@cref}{{[originallemma][0][]0FVJ}{[1][0][]source}{}{}{}}
\makeatother
\crefname{originallemma}{original lemma}{original lemmas}
\Crefname{originallemma}{Original lemma}{Original lemmas}
\makeatletter
\newlabel{chow-lemma-pullback-base-change-chern-classes}{{0FVN}{source}{}{sourcepointer-2}{}}
\newlabel{chow-lemma-pullback-base-change-chern-classes@cref}{{[originallemma][0][]0FVN}{[1][0][]source}{}{}{}}
\makeatother
\crefname{originallemma}{original lemma}{original lemmas}
\Crefname{originallemma}{Original lemma}{Original lemmas}
\makeatletter
\newlabel{chow-lemma-splitting-principle}{{02UL}{source}{}{sourcepointer-3}{}}
\newlabel{chow-lemma-splitting-principle@cref}{{[originallemma][0][]02UL}{[1][0][]source}{}{}{}}
\makeatother
\crefname{originalremark}{original remark}{original remarks}
\Crefname{originalremark}{Original remark}{Original remarks}
\makeatletter
\newlabel{chow-remark-the-proof-shows-more}{{0FVE}{source}{}{sourcepointer-4}{}}
\newlabel{chow-remark-the-proof-shows-more@cref}{{[originalremark][0][]0FVE}{[1][0][]source}{}{}{}}
\makeatother
\crefname{originalsection}{original section}{original sections}
\Crefname{originalsection}{Original section}{Original sections}
\makeatletter
\newlabel{chow-section-change-base}{{0FVF}{source}{}{sourcepointer-5}{}}
\newlabel{chow-section-change-base@cref}{{[originalsection][0][]0FVF}{[1][0][]source}{}{}{}}
\makeatother
\crefname{originalreference}{original reference}{original references}
\Crefname{originalreference}{Original reference}{Original references}
\makeatletter
\newlabel{lemma-cycle-classes}{{0FIF}{source}{}{sourcepointer-6}{}}
\newlabel{lemma-cycle-classes@cref}{{[originalreference][0][]0FIF}{[1][0][]source}{}{}{}}
\makeatother
\crefname{originalreference}{original reference}{original references}
\Crefname{originalreference}{Original reference}{Original references}
\makeatletter
\newlabel{remark-trace}{{0FIK}{source}{}{sourcepointer-7}{}}
\newlabel{remark-trace@cref}{{[originalreference][0][]0FIK}{[1][0][]source}{}{}{}}
\makeatother
\title{Verifying the axioms after a field extension}
\author{The Stacks Project Authors}
\date{stacks-project:weil.tex:local-snapshot-20261012; GNU FDL 1.2 or later}
\begin{document}
\maketitle
\noindent\textit{Editorial note.} All delivered mathematical text is the author's own, copied from the pinned Stacks Project chapter file weil.tex; no statement or proof step was rewritten, reordered or removed.\par\medskip
\noindent\textit{Editorial note.} The counted proof is the author's complete argument: the linear algebra principles transferring the formal axioms, the compatibility of the cycle class map with base change through Chern characters and the splitting principle, and the treatment of axioms (A6) and (A9) through the Kunneth decomposition of the diagonal and the uniqueness of the degree map.\par\medskip
\noindent\textit{Editorial note.} The label lemma-check-over-extension names the counted statement, and the proof environment immediately following it in the source is delivered verbatim as its proof.\par\medskip
\noindent\textit{Editorial note.} The counted statement is carried by the author's own lemma environment in the chapter source.\par\medskip
\noindent\textit{Editorial note.} Pre-open gate: python operations/agent-authoring/tools/precard\_gate.py --source weil.tex --label lemma-check-over-extension --start 5444 --end 5616 --local-id 28125 returned RESULT=OK (L28125 free, occupancy FREE, public interval clear). Pre-assembly tag self-check: canonical tag resolved and not present in reference/public-items.jsonl.\par\medskip
\noindent\textit{Editorial note.} The statement and proof are copied from the pinned chapter file; the cited inputs are kept as external pointers with their canonical Stacks tags.\par\medskip
\noindent\textit{Allowed standard prerequisites (editorial).}\begin{itemize}
\item The cycle class map as the Chern character of a K theory class. Same chapter, the construction of the cycle class map from the Chern character using axioms (A1) to (A4). Expresses the cycle class map as the Chern character so that compatibility with base change can be checked on Chern classes of vector bundles.
\item Compatibility of the base change map on Chow groups with capping by Chern classes and the splitting principle. Chow Homology, the lemma on the pullback along a base change of the base scheme, its compatibility with Chern classes, and the splitting principle for vector bundles on projective space bundles. Reduces the commutation of the cycle class map with base change to the agreement of first Chern classes, which is one of the hypotheses.
\item Uniqueness of the degree map determined by axiom (A6). Same chapter, the remark on the trace map, and the construction in the proof of axiom (A6). Identifies the degree maps after base change, which is what makes axiom (A9) follow by a diagram chase.
\end{itemize}\medskip
\ifdefined\hypertarget\hypertarget{sourcepointer-1}{}\fi
\noindent\textit{Source reference.} Original lemma 0FVJ: sources/stacks-chow/tags.txt, source lines 19814--19814. Prerequisite pointer invoked at proof line(s) 5528; it names a standard fact of a sibling chapter that the counted argument uses, and every step of the counted proof built on it is retained verbatim in the delivered body.\par\medskip
\ifdefined\hypertarget\hypertarget{sourcepointer-2}{}\fi
\noindent\textit{Source reference.} Original lemma 0FVN: sources/stacks-chow/tags.txt, source lines 19818--19818. Prerequisite pointer invoked at proof line(s) 5535; it names a standard fact of a sibling chapter that the counted argument uses, and every step of the counted proof built on it is retained verbatim in the delivered body.\par\medskip
\ifdefined\hypertarget\hypertarget{sourcepointer-3}{}\fi
\noindent\textit{Source reference.} Original lemma 02UL: sources/stacks-chow/tags.txt, source lines 3736--3736. Prerequisite pointer invoked at proof line(s) 5548; it names a standard fact of a sibling chapter that the counted argument uses, and every step of the counted proof built on it is retained verbatim in the delivered body.\par\medskip
\ifdefined\hypertarget\hypertarget{sourcepointer-4}{}\fi
\noindent\textit{Source reference.} Original remark 0FVE: sources/stacks-chow/tags.txt, source lines 19809--19809. Prerequisite pointer invoked at proof line(s) 5549; it names a standard fact of a sibling chapter that the counted argument uses, and every step of the counted proof built on it is retained verbatim in the delivered body.\par\medskip
\ifdefined\hypertarget\hypertarget{sourcepointer-5}{}\fi
\noindent\textit{Source reference.} Original section 0FVF: sources/stacks-chow/tags.txt, source lines 19810--19810. Prerequisite pointer invoked at proof line(s) 5525; it names a standard fact of a sibling chapter that the counted argument uses, and every step of the counted proof built on it is retained verbatim in the delivered body.\par\medskip
\ifdefined\hypertarget\hypertarget{sourcepointer-6}{}\fi
\noindent\textit{Source reference.} Original reference 0FIF: sources/stacks-weil/pinned-original-source.tex, source lines 4713--4713. Prerequisite pointer invoked at proof line(s) 5523; it names a standard fact of the same chapter that the counted argument uses, and every step of the counted proof built on it is retained verbatim in the delivered body.\par\medskip
\ifdefined\hypertarget\hypertarget{sourcepointer-7}{}\fi
\noindent\textit{Source reference.} Original reference 0FIK: sources/stacks-weil/pinned-original-source.tex, source lines 4929--4929. Prerequisite pointer invoked at proof line(s) 5606; it names a standard fact of the same chapter that the counted argument uses, and every step of the counted proof built on it is retained verbatim in the delivered body.\par\medskip
\begin{lemma}
\label{lemma-check-over-extension}
Let $k'/k$ be an extension of fields. Let $F'/F$ be an extension
of fields of characteristic $0$. Assume given
\begin{enumerate}
\item data (D0), (D1), (D2') for $k$ and $F$ denoted
$F(1), H^*, c_1^H$,
\item data (D0), (D1), (D2') for $k'$ and $F'$ denoted
$F'(1), (H')^*, c_1^{H'}$, and
\item an isomorphism $F(1) \otimes_F F' \to F'(1)$, functorial isomorphisms
$H^*(X) \otimes_F F' \to (H')^*(X_{k'})$ on the category of smooth projective
schemes $X$ over $k$ such that the diagrams
$$
\xymatrix{
\Pic(X) \ar[r]_{c_1^H} \ar[d] & H^2(X)(1) \ar[d] \\
\Pic(X_{k'}) \ar[r]^{c_1^{H'}} & (H')^2(X_{k'})(1)
}
$$
commute.
\end{enumerate}
In this case, if $F'(1), (H')^*, c_1^{H'}$ satisfy axioms (A1) -- (A9),
then the same is true for $F(1), H^*, c_1^H$.
\end{lemma}

\begin{proof}
We go by the axioms one by one.

\medskip\noindent
Axiom (A1). We have to show $H^*(\emptyset) = 0$ and that
$(i^*, j^*) : H^*(X \amalg Y) \to H^*(X) \times H^*(Y)$
is an isomorphism where $i$ and $j$ are the coprojections.
By the functorial nature of the isomorphisms
$H^*(X) \otimes_F F' \to (H')^*(X_{k'})$ this
follows from linear algebra: if $\varphi : V \to W$ is an $F$-linear map
of $F$-vector spaces, then $\varphi$ is an isomorphism if and only if
$\varphi_{F'} : V \otimes_F F' \to W \otimes_F F'$ is an isomorphism.

\medskip\noindent
Axiom (A2). This means that given a morphism $f : X \to Y$ of smooth projective
schemes over $k$ and an invertible $\mathcal{O}_Y$-module $\mathcal{N}$
we have $f^*c_1^H(\mathcal{L}) = c_1^H(f^*\mathcal{L})$. This is immediately
clear from the corresponding property for $c_1^{H'}$, the commutative
diagrams in the lemma, and the fact that the canonical map
$V \to V \otimes_F F'$ is injective for any $F$-vector space $V$.

\medskip\noindent
Axiom (A3). This follows from the principle stated in the proof of
axiom (A1) and compatibility of $c_1^H$ and $c_1^{H'}$.

\medskip\noindent
Axiom (A4). Let $i : Y \to X$ be the inclusion of an effective
Cartier divisor over $k$ with both $X$ and $Y$ smooth and projective
over $k$. For $a \in H^*(X)$ with
$i^*a = 0$ we have to show $a \cup c_1^H(\mathcal{O}_X(Y)) = 0$.
Denote $a' \in (H')^*(X_{k'})$ the image of $a$.
The assumption implies that $(i')^*a' = 0$ where $i' : Y_{k'} \to X_{k'}$
is the base change of $i$. Hence we get
$a' \cup c_1^{H'}(\mathcal{O}_{X_{k'}}(Y_{k'})) = 0$ by the axiom
for $(H')^*$. Since $a' \cup c_1^{H'}(\mathcal{O}_{X_{k'}}(Y_{k'}))$
is the image of $a \cup c_1^H(\mathcal{O}_X(Y))$ we conclude by
the princple stated in the proof of axiom (A2).

\medskip\noindent
Axiom (A5). This means that $H^*(\Spec(k)) = F$ and that for $X$ and $Y$ smooth
projective over $k$ the map $H^*(X) \otimes_F H^*(Y) \to H^*(X \times Y)$,
$a \otimes b \mapsto p^*(a) \cup q^*(b)$ is an isomorphism
where $p$ and $q$ are the projections. This follows from the principle
stated in the proof of axiom (A1).

\medskip\noindent
We interrupt the flow of the arguments to show that for every
smooth projective scheme $X$ over $k$ the diagram
$$
\xymatrix{
\CH^*(X) \ar[r]_-\gamma \ar[d]_{g^*} & \bigoplus H^{2i}(X)(i) \ar[d] \\
\CH^*(X_{k'}) \ar[r]^-{\gamma'} & \bigoplus (H')^{2i}(X_{k'})(i)
}
$$
commutes. Observe that we have $\gamma$ as we know axioms
(A1) -- (A4) already; see Lemma \ref{lemma-cycle-classes}.
Also, the left vertical arrow is the one discussed in
Chow Homology, Section \ref{chow-section-change-base}
for the morphism of base schemes $g : \Spec(k') \to \Spec(k)$.
More precisely, it is the map given in
Chow Homology, Lemma \ref{chow-lemma-pullback-base-change}.
Pick $\alpha \in \CH^*(X)$. Write $\alpha = ch(\beta) \cap [X]$
in $\CH^*(X) \otimes \mathbf{Q}$
for some $\beta \in K_0(\textit{Vect}(X)) \otimes \mathbf{Q}$
so that $\gamma(\alpha) = ch^{H}(\beta)$; this is our construction of $\gamma$.
Since the map of Chow Homology, Lemma \ref{chow-lemma-pullback-base-change}
is compatible with capping with Chern classes by
Chow Homology, Lemma \ref{chow-lemma-pullback-base-change-chern-classes}
we see that $g^*\alpha = ch((X_{k'} \to X)^*\beta) \cap [X_{k'}]$.
Hence $\gamma'(g^*\alpha) = ch^{H'}((X_{k'} \to X)^*\beta)$.
Thus commutativity of the diagram will hold if for any locally
free $\mathcal{O}_X$-module $\mathcal{E}$ of rank $r$ and $0 \leq i \leq r$
the element $c_i^H(\mathcal{E})$ of $H^{2i}(X)(i)$
maps to the element $c_i^{H'}(\mathcal{E}_{k'})$ in $(H')^{2i}(X_{k'})(i)$.
Because we have the projective space bundle formula for both
$X$ and $X'$ we may replace $X$ by a projective space bundle
over $X$ finitely many times to show this. Thus we may assume
$\mathcal{E}$ has a filtration whose graded pieces are
invertible $\mathcal{O}_X$-modules
$\mathcal{L}_1, \ldots, \mathcal{L}_r$.
See Chow Homology, Lemma \ref{chow-lemma-splitting-principle} and
Remark \ref{chow-remark-the-proof-shows-more}.
Then $c^H_i(\mathcal{E}$ is the $i$th elementary symmetric polynomial
in $c^H_1(\mathcal{L}_1), \ldots, c^H_1(\mathcal{L}_r)$
and we conclude by our assumption that we have agreement for
first Chern classes.

\medskip\noindent
Axiom (A6). Suppose given $F$-vector spaces
$V$, $W$, an element $v \in V$, and a tensor $\xi \in V \otimes_F W$.
Denote $V' = V \otimes_F F'$, $W' = W \otimes_F F'$ and $v'$, $\xi'$
the images of $v$, $\xi$ in $V'$, $V' \otimes_{F'} W'$. The linear algebra
principle we will use in the proof of axiom (A6) is the following:
there exists an $F$-linear map $\lambda : W \to F$ such that
$(1 \otimes \lambda)\xi = v$ if and only if there exists an $F'$-linear
map $\lambda' : W \otimes_F F' \to F'$ such that
$(1 \otimes \lambda')\xi' = v'$.

\medskip\noindent
Let $X$ be a nonempty equidimensional smooth projective scheme
over $k$ of dimension $d$. Denote $\gamma = \gamma([\Delta])$
in $H^{2d}(X \times X)(d)$ (unadorned fibre products will be over $k$).
Observe/recall that this makes sense as we know axioms (A1) -- (A4) already;
see Lemma \ref{lemma-cycle-classes}. We may decompose
$$
\gamma = \sum \gamma_i, \quad
\gamma_i \in H^i(X) \otimes_F H^{2d - i}(X)(d)
$$
in the K\"unneth decomposition. Similarly, denote
$\gamma' = \gamma([\Delta']) = \sum \gamma'_i$
in $(H')^{2d}(X_{k'} \times_{k'} X_{k'})(d)$.
By the linear algebra princple mentioned above, it suffices
to show that $\gamma_0$ maps to $\gamma'_0$ in
$(H')^0(X) \otimes_{F'} (H')^{2d}(X')(d)$.
By the compatibility of K\"unneth decompositions
we see that it suffice to show that $\gamma$ maps to
$\gamma'$ in
$$
(H')^{2d}(X_{k'} \times_{k'} X_{k'})(d) = (H')^{2d}((X \times X)_{k'})(d)
$$
Since $\Delta_{k'} = \Delta'$ this follows from the discussion above.

\medskip\noindent
Axiom (A7). This follows from the linear algebra fact: a
linear map $V \to W$ of $F$-vector spaces is injective
if and only if $V \otimes_F F' \to W \otimes_F F'$ is injective.

\medskip\noindent
Axiom (A8). Follows from the linear algebra fact used in
the proof of axiom (A1).

\medskip\noindent
Axiom (A9). Let $X$ be a nonempty smooth projective scheme over $k$
equidimensional of dimension $d$. Let $i : Y \to X$ be a nonempty
effective Cartier divisor smooth over $k$.
Let $\lambda_Y$ and $\lambda_X$ be as in axiom (A6) for $X$ and $Y$.
We have to show: for $a \in H^{2d - 2}(X)(d - 1)$
we have $\lambda_Y(i^*(a)) = \lambda_X(a \cup c_1^H(\mathcal{O}_X(Y))$.
By Remark \ref{remark-trace}
we know that $\lambda_X : H^{2d}(X)(d) \to F$ and
$\lambda_Y : H^{2d - 2}(Y)(d - 1)$ are uniquely
determined by the requirement in axiom (A6).
Having said this, it follows from our proof of axiom (A6) for $H^*$ above
that $\lambda_X \otimes \text{id}_{F'}$ corresponds to $\lambda_{X_{k'}}$
via the given identification $H^{2d}(X)(d) \otimes_F F' = H^{2d}(X_{k'})(d)$.
Thus the fact that we know axiom (A9) for $F'(1), (H')^*, c_1^{H'}$
implies the axiom for $F(1), H^*, c_1^H$ by a diagram chase.
This completes the proof of the lemma.
\end{proof}

\end{document}
